\subsection{忠实平坦环与有限性条件}

命题 11. 设 B 是环，A 是 B 的子环，且 B 是忠实平坦右 A-模。为使左 A-模 F 是有限生成的（相应地有限表示的），必须且只需 B-模 \(B \otimes_{A} F\) 是有限生成的（相应地有限表示的）。

\begin{enumerate}
\def\labelenumi{(\arabic{enumi})}
\item
  不对 B 作任何假设，显然，若 F 是有限生成的左 A-模，则 \(B \otimes_{A} F\) 是有限生成的左 B-模。反之，若 \(B \otimes_{A} F\) 是有限生成的 B-模，则它由形如 \(1 \otimes x_{i}\) （其中 \(x_{i} \in F\) ）的有限个元素生成；若 M 是由诸 \(x_{i}\) 生成的 F 的子 A-模，j 是典范单射 \(M \to F\) ，则 \(1_{B} \otimes j: B \otimes_{A} M \to B \otimes_{A} F\) 是满同态，于是 j 是满射（第 1 目，命题 2），这就证明 F 是有限生成的。
\item
  若 \(\mathbf{F}\) 有有限表示，则 \(\mathbf{B} \otimes_{\mathbf{A}} \mathbf{F}\) 亦然，这不需要对 B 作任何假设（§2 第 8 目）。剩下要证明：若 \(B \otimes_{A} F\) 有有限表示，则 F 亦然。我们已由 (1) 知道 F 是有限生成的，因此存在满同态 u: L → F，其中 L 是有限生成的自由 A-模。设 R 是 u 的核，则 \(B \otimes_{A} R\) 等同于满同态 \(1_{B} \otimes u: B \otimes_{A} L \to B \otimes_{A} F\) 的核（§2 第 3 目，注 2）。由于由假设 \(B \otimes_{A} F\) 有有限表示，我们断言（§2 第 8 目，引理 9） \(B \otimes_{A} R\) 是有限生成的；于是由 (1) 推出 R 是有限生成的 A-模，从而 F 有有限表示。
\end{enumerate}

命题 12. 设 B 是环，A 是 B 的中心中的交换子环，且 B 是忠实平坦 A-模。为使 A-模 F 是投射的且有限生成的，必须且只需 \(B \otimes_{A} F\) 是有限生成的投射左 B-模。

不对 A 或 B 作任何假设，这个条件显然是必要的（《代数学》第二章 §5 第 1 目命题 4 的推论）；我们证明它是充分的。由于有限生成的投射模有有限表示（§2 第 8 目，引理 8），假设蕴含 F 依命题 11 有有限表示，因此，对每个 A-模 M，有典范同构

\[
\omega : \mathbf {B} \otimes_ {\mathbf {A}} \operatorname{Hom} _ {\mathbf {A}} (\mathbf {F}, \mathbf {M}) \rightarrow \operatorname{Hom} _ {\mathbf {B}} (\mathbf {B} \otimes_ {\mathbf {A}} \mathbf {F}, \mathbf {B} \otimes_ {\mathbf {A}} \mathbf {M})
\]

（§2 第 10 目，命题 11）。然后设 \(v: \mathbf{M} \to \mathbf{M}''\) 是满的 A-模同态，并考虑交换图

\[
\begin{array}{c} \mathrm{B} \otimes_ {\mathrm{A}} \operatorname{Hom} _ {\mathrm{A}} (\mathrm{F}, \mathrm{M}) \xrightarrow {\omega} \operatorname{Hom} _ {\mathrm{B}} (\mathrm{B} \otimes_ {\mathrm{A}} \mathrm{F}, \mathrm{B} \otimes_ {\mathrm{A}} \mathrm{M}) \\ 1 _ {\mathrm{B}} \otimes \operatorname{Hom} (1 _ {\mathrm{F}}, v) \Bigg \downarrow \quad \Bigg \downarrow \operatorname{Hom} (1 _ {\mathrm{B}} \otimes \mathrm{F}, 1 _ {\mathrm{B}} \otimes v) \\ \mathrm{B} \otimes_ {\mathrm{A}} \operatorname{Hom} _ {\mathrm{A}} (\mathrm{F}, \mathrm{M} ^ {\prime \prime}) \xrightarrow [ \omega ]{\longrightarrow} \operatorname{Hom} _ {\mathrm{B}} (\mathrm{B} \otimes_ {\mathrm{A}} \mathrm{F}, \mathrm{B} \otimes_ {\mathrm{A}} \mathrm{M} ^ {\prime \prime}) \end{array}
\]

由于 \(1_{B} \otimes v\) 是满射且 \(B \otimes_{A} F\) 假设为投射的， \(\operatorname{Hom}(1_{B} \otimes F, 1_{B} \otimes v)\) 是满射（《代数学》第二章 §2 第 2 目，命题 4），因此 \(1_{B} \otimes \operatorname{Hom}(1_{F}, v)\) 也是满射。但由于 B 是忠实平坦 A-模， \(\operatorname{Hom}(1_{F}, v)\) 本身是满射（第 1 目，命题 2），因此 F 是投射 A-模（《代数学》第二章 §2 第 2 目，命题 4）。

